\writeSectionFrame{\presentationTitle}{Fazit}
\begin{frame}{Einsetzen wenn ...}
	\begin{itemize}
 		\item Algorithmen austauschbar sein sollen.
        \item es Varianten von Algorithmen geben soll.
 		\item Algorithmen und algorithmusspezifischen Datenstrukturen gekapselt werden sollen.
 		\item die Auswahl mehrerer Algorithmen möglich sein soll
        \item unterschiedliche Verhaltensweisen in eigene Klassen extrahiert werden können.
 	\end{itemize}
\end{frame}

\note{
	 Damit das Strategiemuster eingesetzt werden kann, muss ein Teil der Aufgabenstellung in einen austauschbaren Algorithmus extrahierbar sein. Algorithmus kann auch eine bestimmte Vorgehensweise sein, z.B. ein Dokument wird als pdf, txt oder csv-Dokument dargestellt. Der Algorithmus wird vollständig gekapselt. Er kann also auf Datenstrukturen u.ä. zugreifen, ohne dass das restliche Programm etwas davon weiß. Und ein wichtiger Aspekt für den Einsatz dieses Musters ist die Austauschbarkeit des Algorithmus.
}

\begin{frame}{Einsatzbeispiele}
	\begin{itemize}
 		\item Sortieralgorithmen
 		\item Validationsregeln
 		\item Textformatierung
        \item Datenbankzugriffe
        \item Wegeberechnung
        \item Testverfahren
        \item Logging
 	\end{itemize}
\end{frame}

\realworld{Java Collections Framework}{Sortieralgorithmen}
\note{
    Im Java-Collections-Framework wird das Strategie-Muster verwendet, um verschiedene Sortieralgorithmen anzuwenden. Die Collections.sort()-Methode kann eine benutzerdefinierte Comparator-Implementierung akzeptieren, die als Strategie fuer das Sortieren von Objekten verwendet wird.
}

\realworld{Java Stream API}{filter(), map(), reduce() usw.}
\note{
   Die Stream-API von Java verwendet das Strategie-Muster für Operationen wie filter, map, und reduce. Diese Operationen akzeptieren funktionale Interfaces, die als Strategien für die Verarbeitung von Daten verwendet werden.
}

\realworld{Spring Framework}{Datenbankkonnektivität und -zugriff}
\note{
    Im Spring Framework können verschiedene Strategien für die Datenbankkonnektivität und -zugriff verwendet werden. Zum Beispiel kann Spring Data verschiedene Repository-Implementierungen (wie JPA, MongoDB, etc.) verwenden, je nach den Konfigurationsparametern.
}

\realworld{Apache Commons CLI}{Kommandozeilen-Optionen}
\note{
    Die Apache Commons CLI-Bibliothek verwendet das Strategie-Muster für die Verarbeitung von Kommandozeilen-Optionen. Sie ermöglicht es, verschiedene Parser-Strategien zu definieren, um unterschiedliche Arten von Kommandozeilenargumenten zu verarbeiten.
}

\realworld{Log4J}{Log-Levels und Formatierungen}
\note{
   In Log4j werden verschiedene Strategien für das Logging verwendet, z. B. unterschiedliche Log-Level (DEBUG, INFO, WARN, ERROR) und Layouts (wie PatternLayout oder XMLLayout).
}

\vorteil{Austauschbarkeit der Algorithmen während Laufzeit}
\note{
	Das Muster ermöglicht es, die Implementierung einer Strategie zur Laufzeit zu ändern, ohne den Kontext oder andere Teile des Systems zu beeinflussen.
}

\vorteil{Erweiterbarkeit}
\note{
	Neue Strategien können einfach hinzugefügt werden, ohne den bestehenden Code zu ändern. Dies fördert eine hohe Erweiterbarkeit und Wartbarkeit des Systems.
}

\vorteil{Trennung von Verantwortlichkeiten}
\note{
	Die Logik der verschiedenen Algorithmen ist vom Kontext (der Klasse, die die Strategie verwendet) getrennt. Dies führt zu einem klareren, modulareren Design.
}

\vorteil{Erleichterte Wartung}
\note{
	Da der Kontext und die Strategien getrennt sind, können Änderungen an einer Strategie ohne Beeinflussung des Kontexts vorgenommen werden.
}

\vorteil{Wiederverwendbarkeit}
\note{
	Strategien können in verschiedenen Kontexten wiederverwendet werden, wodurch der Code wiederverwendbar und effizienter wird.
}

\vorteil{Einfachere Tests}
\note{
	Durch die Trennung der Strategien können diese unabhängig getestet werden. Der Kontext kann ebenfalls isoliert getestet werden, indem man Mock- oder Stub-Strategien verwendet.
}

\vorteil{Open-Closed-Prinzip}
\vorteil{Single-Responsibility-Prinzip}
\note{
	Diese Vorteile kann man im Grunde alle darauf herunterbrechen, dass Strategie das Open-Closed und das Single-Responsibility-Prinzip einhält. 
}

\nachteil{Verwaltung der Strategie-Instanzen}
\note{
	Es kann zusätzliche Logik erforden, um die Strategie-Instanzen zu verwalten und auszuwählen, insbesondere wenn viele verschiedene Strategien existieren.
}

\nachteil{Potentielle Leistungseinbußen}
\note{
	Die Nutzung von Strategien zur Laufzeit kann zu einem gewissen Overhead führen, insbesondere wenn häufig zwischen Strategien gewechselt wird. Eventuell müssen bei größeren Systemen Gegenmaßnahmen wie z.B. Caching oder Einschränkungen in der Strategie-Auswahl vorgenommen werden. Dynamisches Binden kann zusätzliche Laufzeit erfordern. 
}

\nachteil{Zusätzliche Indirektionsebene}
\note{
	Wenn eine Methode in der Strategie aufgerufen wird, geschieht dies in der Regel durch eine Referenz auf das Strategy-Interface oder eine Basisklasse. Dies erfordert eine zusätzliche Ebene der Indirektion, da der Aufruf der Methode über die Basisklasse oder das Interface erfolgt, anstatt direkt auf die konkrete Implementierung zuzugreifen. Diese zusätzliche Ebene kann zu einem gewissen Laufzeit-Overhead führen.
}

\nachteil{Eingeschränkte Sichtbarkeit der Strategie-Auswahl}
\note{
	In komplexen Systemen kann es schwierig sein, nachzuvollziehen, welche Strategie zur Laufzeit gewählt wird, was zu einer verringerten Klarheit führen kann.
}